\SourcePage{GELFAND-P000463}{unlabelled}
Consid.
\[
K_*[x]=\operatorname{Hom}_K(K[x],K).
\]
(“diff. operators”).
\[
(x_*^n,x^m)=\delta_{n,m}
\]
Multiplication:
\[
\underline{x_*^n\cdot x_*^m=C_{n+m}^n x_*^{n+m}}
\qquad\text{— \underline{table}.}
\]
Here $K$ plays no role.

It suffices to take $K=\mathbb Z_p$ or $\mathbb Q$.
\[
K_*[x]=\prod_{q=1}^\infty B_q,\qquad
B_q=\mathbb Z_p[y_{p^i}]/(y_{p^i}^{\,p})
\]
\Unclear{GELFAND-P000463}{product-factorization}{The left side reads $K_*[x]$. The product has lower limit $q=1$ and upper limit $\infty$.}


$y_1,\ldots,y_i,\ldots,y_{p^i}$ — generators.
\Unclear{GELFAND-P000463}{generator-powers}{In three places the $i$ next to $p$ is raised. This is set as $y_{p^i}$. In the denominator the separate outer exponent $p$ on $y$ is visible.}



\Needspace{7\baselineskip}

\SourcePage{GELFAND-P000464}{unlabelled}
\underline{30/\mbox{$\overline{\underline{\mathrm X}}$}-72.}
\quad\underline{Gelfand sem..}

\underline{I. Bernstein}

\uuline{И.М.Г.}

\SourceFigure{GELFAND-P000464}{general-graph}{The two arrows at left point to the same vertex. From the middle vertex arrows lead left, right and down to the right. The bent line at right has arrowheads. The small strokes at vertices have no established reading as letters.
\par\smallskip\centering\color{black}
[DIAGRAM]}


There is a directed graph. This is a category.

\sout{Р}\quad Consid. some diagrams, with at the points — vect. spaces,
on the arrows maps (functoriality is not required).
\Unclear{GELFAND-P000464}{cancel-before-consider}{A loop survives at the start of the crossed-out sign. Possibly Cyrillic «Р».}

\Unclear{GELFAND-P000464}{diagram-adjective}{The words before “diagrams” provisionally read «какие л.» (“some”).}


Prob. — to classify.

\SourceFigure{GELFAND-P000464}{one-arrow}{A slightly curved rightward arrow joins the two points.
\par\smallskip\centering\color{black}
[DIAGRAM]}


Prob. of classif. of maps.\quad The only invs. are $\dim\operatorname{Ker}$ and $\dim\operatorname{Coker}$.
Reduces to an embedding.

\SourceFigure{GELFAND-P000464}{one-loop}{One loop. Its arrowhead on the left side points upward.
\par\smallskip\centering\color{black}
[DIAGRAM]}


one transformation.

\SourceFigure{GELFAND-P000464}{two-loop}{Two loops share the left vertex. The top of the outer loop is faint. Overlapping arrowhead strokes are visible near the vertex.
\par\smallskip\centering\color{black}
[DIAGRAM]}


— the prob. is unsolvable (a pair of arbitrary lin. operators). \sout{\ldots}
\Unclear{GELFAND-P000464}{cancel-before-ona}{The short entry before “it” is crossed out and illegible.}


It reduces to the prob. of commuting ones (Gelf. and Ponomarev).
\Unclear{GELFAND-P000464}{wild-name}{The first letter is overwritten. The abbreviation reads «Гельф.» (Gelf.).}


\sout{Reduc.\ldots}
\Unclear{GELFAND-P000464}{large-cancel-fragment}{The crossed-out line may begin «Свод.» (“reduc.”).}

\Incomplete{GELFAND-P000464}{large-cancel-remainder}{The remainder of the crossed-out line is illegible.}


\[
\left.
\begin{gathered}
\text{Space and subspace.}\\
\text{Space and 2 subspaces}
\end{gathered}
\right\}\quad\text{triv.}
\]
Can be done:\quad Space and \uwave{\mbox{$3$} lin. subspaces}.
\Unclear{GELFAND-P000464}{three-space-word}{The start of the line provisionally reads «Делается» (“can be done”).}


There are $9$, invariants, not $8$!

\SourceFigure{GELFAND-P000464}{three-inward}{Three arrows point toward the upper vertex. Their arrowheads partly merge.
\par\smallskip\centering\color{black}
[DIAGRAM]}

(problem).

\SourceFigure{GELFAND-P000464}{four-inward}{Four edges meet at the heavily inked central vertex. Overlapping strokes obscure the direction of two edges.
\par\smallskip\centering\color{black}
[DIAGRAM]}


Space and $4$ subspaces classified (Gelf. and Ponom.) —
\sout{\ldots} already an inf. number of indecomp.
\Incomplete{GELFAND-P000464}{cancel-after-names}{The word after the dash is crossed out and illegible.}


\underline{Gabriel}: in which graphs are there finitely many indecomp. systems?

Answer: draw a Dynkin diagram! — these and only these are such graphs.

Exceptions: diagrams with double (generally multiple) arrows.


\Needspace{7\baselineskip}

\SourcePage{GELFAND-P000465}{unlabelled}
\underline{Theorem}\quad
(\begin{tabular}[t]{@{}c@{}}Roytov\\[-.45ex]\scriptsize from Kyiv\end{tabular}
or Nazarovy).
\Unclear{GELFAND-P000465}{theorem-names}{The names provisionally read «Ройтов» and «Назаровы». The addition below reads “from Kyiv.”}


A class of diagrams is descr. with a classific. as in Gelf. and Ponomarev.

\uwave{Connection \sout{\ldots}\ with singularities!}\quad — see Artin, Tyurina, Brieskorn.
\Incomplete{GELFAND-P000465}{cancel-specialities}{The word between “Connection” and “with singularities” is crossed out and illegible.}


\begin{center}
\rule[.5ex]{18mm}{.35pt}\enspace$\cdot$\enspace\rule[.5ex]{25mm}{.35pt}
\end{center}

\uuline{И.Б.}\quad \uwave{Category of} vect. spaces over points and maps over arrows.

(The graph is fixed). \quad \uwave{Def. direct sum}. This is $\operatorname{Cat}(\Gamma)$
($\Gamma$ — graph)

Morphisms — maps over each point with commut. diagrs.

$\bullet$\qquad category of linear spaces.

Important: no need to ident. isomorphic ones.
\Unclear{GELFAND-P000465}{important-line}{The abbreviation before “isomorphic” provisionally reads «отожд.» (“ident.”).}


It is easy to see that $\forall$ object $=\sum$ indecomp. ones. (uniquely up to isom.).

$1)$\quad \uwave{Classif. \mbox{$\Gamma$}, with finitely many indecomp. objects.}

\underline{Restrict to connected ones}:

\SourceFigure{GELFAND-P000465}{dynkin-a-mark}{A small angular stroke rises above the third vertex from the left.
\par\smallskip\centering\color{black}
[DIAGRAM]}


\SourceFigure{GELFAND-P000465}{dynkin-d}{There is a small leftward arrowhead on the segment between the third and fourth vertices.
\par\smallskip\centering\color{black}
[DIAGRAM]}


\SourceFigure{GELFAND-P000465}{dynkin-e}{There are six vertices on the horizontal chain. A branch descends from the third vertex from the left.
\par\smallskip\centering\color{black}
[DIAGRAM]}

$k=6,7,8$\qquad (necessarily $2$ segments on the left).

Place the arrows arbitrarily!!

$D_4$ — space and $3$ subspaces.

Idea: no cycles (one obt. eigenvalues), no
\Unclear{GELFAND-P000465}{idea-word}{The abbreviation in parentheses may read «получ.» (“obt.”).}


\SourceFigure{GELFAND-P000465}{forbidden-four}{Four segments meet at the black dot. No separate vertices are visible at the ends.
\par\smallskip\centering\color{black}
[DIAGRAM]}

— there will be an inf. number.

\SourceFigure{GELFAND-P000465}{forbidden-forks}{The chain has a fork at each end. Five vertices are marked on the chain itself.
\par\smallskip\centering\color{black}
[DIAGRAM]}

— also imposs., since ident.
\Unclear{GELFAND-P000465}{fork-explanation}{The last abbreviation provisionally reads «отожд.» (“ident.”).}


$2)$\quad \uwave{The no. of indecomp. objs. does not depend on the dir. of the arrows.}

\uwave{It is characterized by the set of dimensions}!!

one should consid. the root syst.

$\gamma$ — pos. root,\qquad $\alpha_i$ — simp. roots.

To each root — an indecomp. object
\[
\gamma=\sum n_i\alpha_i,\qquad n_i\text{ — dimensions!!}
\]
\Unclear{GELFAND-P000465}{root-symbol}{Both signs have a lower loop and two upper strokes. They are provisionally set as $\gamma$.}


No. $=$ no. of simple roots.
\Unclear{GELFAND-P000465}{root-count}{The line reads «простых корней» (“simple roots”).}


\SourceFigure{GELFAND-P000465}{a3-example}{Six rows of dimensions stand below the chain with two rightward arrows.
\par\smallskip\centering\color{black}
[DIAGRAM]\par\smallskip
$\begin{array}{ccc}
1&0&0\\
0&1&0\\
0&0&1\\
1&1&0\\
0&1&1\\
1&1&1
\end{array}$}

\Unclear{GELFAND-P000465}{array-overwrite}{The middle 1 in row two is overwritten. A remnant of a loop survives to its right.}



\Needspace{7\baselineskip}

\SourcePage{GELFAND-P000466}{unlabelled}
$3)$\quad \uwave{Consid. the \mbox{$\Gamma$}, for which \mbox{$\operatorname{Cat}(\Gamma)$} is solvable.}

\SourceFigure{GELFAND-P000466}{loop-category}{The two loops share a vertex. Inside the upper loop there is a short slanted stroke separate from its arrowhead.
\par\smallskip\centering\color{black}
[DIAGRAM]}


Unsolvable categ. — if there is an embedding $C\longrightarrow\operatorname{Cat}(\Gamma)$,
taking nonisom. objects to nonisom. ones.

They are classif. and correspond to extended Dynkin diagrams!

\SourceFigure{GELFAND-P000466}{three-arrow}{Three arrows leave one vertex, upward, rightward and down to the left.
\par\smallskip\centering\color{black}
[DIAGRAM]}


\SourceFigure{GELFAND-P000466}{affine-a}{Seven points lie on a straight chain. An arc joins the endpoints.
\par\smallskip\centering\color{black}
[DIAGRAM]}


\SourceFigure{GELFAND-P000466}{affine-d}{Seven points lie along the chain, with two branches at each end.
\par\smallskip\centering\color{black}
[DIAGRAM]}


\SourceFigure{GELFAND-P000466}{affine-e}{Seven points lie along the horizontal chain. Two segments with points descend from the third point from the left.
\par\smallskip\centering\color{black}
[DIAGRAM]}


(appear in the theory of \(\overset{\text{\sout{grp.}}}{\vphantom{x}}\) groups generated by reflections.)
\Unclear{GELFAND-P000466}{cancel-in-group-comment}{The crossed-out insertion above the line may read «гр» (“grp.”).}


These are all the graphs for which the prob. is solvable.

\uuline{V. Buchstaber.}

\uuline{Lazur’s thm.}
\Unclear{GELFAND-P000466}{theorem-name}{The name provisionally reads «Лазура» here. The middle vowel may be «а».}

\[
\Lambda=\mathbb Z[a_1,\ldots,a_n,\ldots],\qquad\deg a_n=n\ge1.
\]
\underline{Proof.}
\[
\Lambda\longrightarrow\bar\Lambda=\Lambda/(\text{decomp.})
\]
\Unclear{GELFAND-P000466}{decomposables-quotient}{A horizontal stroke stands above the second $\Lambda$. The denominator reads «разл.» (“decomp.”).}


$\bar\Lambda$ — \uwave{univ. ring for}
\(\underset{\downarrow}{\boxed{\text{\scriptsize such}}}\)
\uwave{polynomials} $h(x,y)$,
\SourceFigure{GELFAND-P000466}{such-insertion}{“Such” is inserted in a bubble before “polynomials.”}


\(\underset{\downarrow}{\boxed{\text{\scriptsize that}}}\)\quad
$h(x,y)=h(y,x)$,\quad $h(x,0)=0$,
\[
h(y,z)+h(x,y+z)=h(x,z)+h(y,x+z).
\]
\SourceFigure{GELFAND-P000466}{that-insertion}{“That” is inserted in a separate bubble before the conditions.}


This is cohomology:\quad $\underset{\text{(cycles)}}{dh}=0$
\qquad (in the bar construction for polynomials)
\Unclear{GELFAND-P000466}{cycle-caption}{Below $dh$ is the word “cycles.”}


\[
\operatorname{Ext}^{*,*}_{A[x]}(A,A)=H^{*,*}(A[x])
\]
\Unclear{GELFAND-P000466}{cohomology-conventions}{The subscript on $\operatorname{Ext}$ reads $A[x]$. There are two stars above the symbol on each side.}

\[
\begin{gathered}
A[x]\longrightarrow A[x]\otimes\bar A[y]\\
x\longrightarrow x+y
\end{gathered}
\]
\Unclear{GELFAND-P000466}{normalized-coproduct}{Below the first arrow the manuscript reads $x\longrightarrow x+y$.}


Cohom. is calc. using resolutions.

The bar constr. is one of the resolutions.

\[
\mathcal A:\quad
A[x]\xrightarrow{\Delta}A[x]\otimes\bar A[y]
\longrightarrow A[x]\otimes\bar A[y]\otimes\bar A[z]\longrightarrow\ldots
\]
$\bar A[y]$ — polynomials with no const. term.

We obt. an acyclic complex which is used to compute cohomology
\[
(\text{put }\underset{h(x,y)}{\text{everywhere}}\text{ — }
\tilde a\Rightarrow a\otimes1+1\otimes a)=\Delta a).
\]
\Unclear{GELFAND-P000466}{primitive-condition}{A wavy stroke stands above $a$. The following sign provisionally reads $\Rightarrow$. Both closing parentheses are visible.}


One must classif. cycles in a particular resolution.
\[
0\longrightarrow B\longrightarrow Z\longrightarrow H^2(A[x])\longrightarrow0
\]
\Unclear{GELFAND-P000466}{exact-sequence}{The term after $B$ is $Z$, and $H$ has superscript $2$.}


$x^n\longrightarrow(x+y)^n-x^n-y^n$\qquad (these are boundaries)


\Needspace{7\baselineskip}

\SourcePage{GELFAND-P000472}{unlabelled}
\SourceFigure{GELFAND-P000472}{top-product-diagram}{The node $C_1$ is heavily retraced. The arrow from $C_i$ to $A$ is dashed, and the arrow from $C_{i+j}$ to $C_j$ has a single break.\par\smallskip\centering
[DIAGRAM]}


We have obtained a homomorphism $C_{i+j}\longrightarrow A_0$.
\Unclear{GELFAND-P000472}{homomorphism-verb}{The first word is read as «Получили» (“we have obtained”). Its last two letters run together.}


This gives
\[
A\underset{\text{assigns}}{\rightsquigarrow}H(A)
\qquad\text{(a ring)}
\]
$H(A)$ is generally noncommutative (!) in a genuine sense.
The classification of multivalued groups rests on this.
\Unclear{GELFAND-P000472}{multivalued-groups}{The manuscript abbreviates “multivalued” as «многозн.».}

(So far in topology only commut. rings have occurred — there one always has Hopf algebras).

\uline{Standard resolution or \mbox{$B$}-resolution.}
\[
\overline A=\sum_{n>0}A_n
\]
\[
C_1=A\underset{A_0}{\otimes}\overline A
\qquad\text{Henceforth }A_0\text{ is omitted}
\]
(The starting rings are commutative!)
\[
\begin{gathered}
C_s=A\otimes(\overline A)^s.\\
d_{s+1}\colon A\otimes(\overline A)^{s+1}
\longrightarrow A\otimes(\overline A)^s
\end{gathered}
\]
\[
\begin{aligned}
d_s\bigl(a_0[a_1|a_2|\ldots|a_s]\bigr)
={}&(-1)^{\varepsilon(0)}a_0a_1[a_2|\ldots a_s]\\
 &+\sum(-1)^{\varepsilon(n)}a_0[a_1|\ldots|a_na_{n+1}|\ldots a_s]
\end{aligned}
\]
\Unclear{GELFAND-P000472}{bar-differential-indices}{Two indices in the product are retraced. The reading is $a_na_{n+1}$.}


$\varepsilon(n)$ — certain functions of the index,
\sout{For even grading \mbox{$\varepsilon(n)=n$}.}
\par\uline{\mbox{$\varepsilon(n)=n$}, if all elems. are even.}
\SourceFigure{GELFAND-P000472}{even-sign-cancellation}{The condition on even elements is added below the deleted line and set off by a rule.}


\[
\operatorname{Hom}_A(C_s,A_0)\simeq(\overline A^{\,*})^s,
\qquad \overline A^{\,*}=\operatorname{Hom}_{A_0}(\overline A,A_0)
\]
\[
d_s^*\colon[\alpha_1|\alpha_2|\ldots|\alpha_{s-1}]
=\sum(-1)^{\varepsilon(n)}
[\alpha_1|\ldots|\overline\Delta\alpha_n|\ldots|\alpha_{s-1}].
\]
\[
\varepsilon(n)=n,\quad\text{if }\deg\alpha_i=2k_i
\]

\uline{Def. \mbox{$\overline\Delta$}.}
\SourceFigure{GELFAND-P000472}{lower-coproduct-box}{The top arrow has a clear right head. Its retraced left end also resembles a head.\par\smallskip\centering
[DIAGRAM]}


In fact
\[
\operatorname{Im}\overline\Delta\subset\overline{A^*}.
\]
\Unclear{GELFAND-P000472}{lower-reduced-image}{The bar on the right extends over both $A$ and the star.}


This gives another definition of multiplication,
\[
\begin{aligned}
C_s^*&\qquad f=[\alpha_1|\ldots|\alpha_s]\\
C_{s'}^*&\qquad g=[\alpha'_1|\ldots|\alpha'_{s'}]
\end{aligned}
\qquad
f\cdot g=[\alpha_1|\ldots|\alpha_s|\alpha'_1|\ldots|\alpha'_{s'}]
\]
\Unclear{GELFAND-P000472}{bottom-prime-index}{The last index in $g$ reads $s^{\prime}$. The separator after the first primed term in the product is heavily retraced.}



\Needspace{7\baselineskip}

\SourcePage{GELFAND-P000473}{unlabelled}
\uline{20/\mbox{$\overline{\underline{\mathrm{XI}}}$}-72.}
\qquad\uline{Gelfand sem..}

\uuline{Kostyuchenko.}\qquad\uuline{One example.}
\SourceFigure{GELFAND-P000473}{date-heading}{The month is given as a Roman numeral. The name and title are double-underlined.}

\[
\begin{gathered}
v''+2a\lambda v'+b\lambda^2v=0\\
v(0)=v(\pi)=0.
\end{gathered}
\qquad\text{in }L_2[0,\pi].
\]

\SourceFigure{GELFAND-P000473}{pde-bracket}{A line under the equation turns up at both ends.\par\smallskip\centering
[DIAGRAM]}

\[
\begin{gathered}
0<x<\pi\\
0\leqslant t<+\infty
\end{gathered}
\]
\Unclear{GELFAND-P000473}{top-time-inequality}{The second inequality sign is retraced. The provisional reading is $0\leqslant t<+\infty$.}

\uline{Suppose the eq. is elliptic.}
\Unclear{GELFAND-P000473}{elliptic-heading}{The first word reads «Пусть» (“Suppose”).}


\[
\underline{a^2<b};\qquad\underline{b>0}.
\]
\[
v_k(x)=e^{\gamma kx}\sin kx,\qquad
\gamma=\frac{-a}{\sqrt{b-a^2}};\qquad k=\pm1,\ \pm2,\ \ldots .
\]
\SourceFigure{GELFAND-P000473}{gamma-bracket}{A rule with upturned ends is drawn below the formula for $\gamma$.}


1. The system is complete, but it is unknown how to find the coefficients.
This is the question of finding the biorthog. system.
\Unclear{GELFAND-P000473}{upper-system-correction}{The end of “system” is retraced. A separate mark is crossed out after the period.}

(found in terms of the inverse fn. of the Christoffel–Schwarz integral)

(\uline{a biorthog. system exists})

2. \uline{This is not a basis!} (not every vector expands into a conv. series in this system)

\sout{Even}\quad\uline{Moreover:} if there is an expansion in $L^2$,
then the limit — an analytic fn.! — and satisfies inf. many boundary conditions (the converse holds too).
\Unclear{GELFAND-P000473}{middle-cancelled-even}{The crossed-out word may read «Даже» (“Even”).}


It cannot be summed by the Abel method of the first order.

\uline{if}
\[
\sum_{\lambda_k>0}e_k e^{\lambda_k t}v_k(x)
\]
\Unclear{GELFAND-P000473}{middle-abel-coefficient}{The coefficient is heavily retraced. The provisional reading is $e_k$.}

\Incomplete{GELFAND-P000473}{middle-abel-cancelled}{Two short entries after $v_k(x)$ are heavily struck out. Their readings are unresolved.}

has a limit as $t\to0$,
then the \uwave{limiting fn.} is analytic.
\SourceFigure{GELFAND-P000473}{limiting-function-line}{A wavy stroke runs above “fn.” and below “limiting”.}


(Discrediting the Fourier method for the elliptic eq.)

With certain nonseparated conditions one can solve it.
\Unclear{GELFAND-P000473}{boundary-abbreviation}{The adjective before “conditions” is abbreviated as «неразд.».}


\Incomplete{GELFAND-P000473}{lower-cancelled-marks}{Two marks before “Reason” are crossed out. Their readings are unresolved.}

\uline{Reason for the poor behav.:} the resolvent grows $\exp$ in the angles
\SourceFigure{GELFAND-P000473}{resolvent-angle-diagram}{Two intersecting lines, broken strokes on the left, a dot at their crossing, and a circle and arrow on the right.\par\smallskip\centering
[DIAGRAM]}


On $[0,\pi-\varepsilon]$ it is overdetermined.

On $[0,\pi+\varepsilon]$, $0<\varepsilon<\pi$ it remains complete !!! (I. Bernstein)
\Unclear{GELFAND-P000473}{epsilon-lower-bound}{The lower bound is heavily retraced and is provisionally set as zero.}


On $[0,2\pi]$ many orthog. fns. arise.

\uline{A problem in the theory of entire fns..}
\[
\int_0^\pi f(x)e^{\gamma kx}\sin kx\,dx=0
\quad\Longrightarrow\quad
F(\lambda)=\int_0^\pi f(x)e^{\gamma\lambda x}\sin\lambda x\,dx
\]
$F(\lambda)$ is an entire fn. of $\exp$ type. It is 0 at integer points
$F(n)=0$, $n>0$.

\uline{The specific features must be taken into account.}

On $[0,\pi-\varepsilon]$ the type and indicator diagram work.

Uniqueness — ?

It is unknown what happens if one replaces $k$ by $\lambda_k$.


\Needspace{7\baselineskip}

\SourcePage{GELFAND-P000474}{unlabelled}
\uuline{I. M. Gelfand}\qquad\uuline{on Arnold's work.}
\SourceFigure{GELFAND-P000474}{gelfand-arnold-heading}{The name and topic are double-underlined.}

\[
(\mathbb R^n,0)\longrightarrow(\mathbb R,0)
\]
$f(x)=0$ — hypersurf.

How does the hypersurf. behave in a neighbourhood of $0$ — ?

$f\in C^\infty$;\quad $f(0)=0$, $df(0)=0$. (crit. point).

(In fact consider $\mathbb C^n\longrightarrow\mathbb C$)

Equiv. hypersurfaces — differ by a change of variables
\[
(\mathbb R^n,0)\longrightarrow(\mathbb R^n,0)\qquad
\text{(diffeom.)}.
\]

\uline{Simple hypersurf.:}\qquad\uline{\mbox{$f(x)=0$}}

if one perturbs it, there will be finitely many orbits
in jets of order $N$ (the bound on the number of orbits does not depend on the index $N$).

Enumeration of the simple orbits:
\SourceFigure{GELFAND-P000474}{simple-chain-graphs}{Two chains are drawn at left. The lower chain forks at its end.\par\smallskip\centering
[DIAGRAM]}

\SourceFigure{GELFAND-P000474}{exceptional-graphs}{The right-hand chain has a downward branch at the third vertex. Its last links are broken.\par\smallskip\centering
[DIAGRAM]}


Next: Most orbits are simple (as a rule, the orbit is simple).
\Incomplete{GELFAND-P000474}{middle-parenthetic-cancelled}{A short fragment after the opening parenthesis is struck out. Its reading is unresolved.}

\uline{Nonsimple ones of codimension 6.}

\[
x_1^{n+1}+x_2^2+\cdots\,\ast=0
\]
— corr. to the first type.
\Unclear{GELFAND-P000474}{middle-simple-equation}{The first letter of the exponent $n+1$ is retraced. A separate star-like mark follows the ellipsis. A rule with upturned ends is drawn below the formula.}


\uline{Versal deformation.}
\[
\begin{gathered}
F(x,\lambda)=0,\qquad F(x,0)=f(x),\\
\lambda=(\lambda_1,\lambda_2,\ldots,\lambda_k)
\end{gathered}
\]
Every other def. is equiv. to a def. induced from this one as follows:
\[
G(x,\mu)=F(X(x,\mu),\Psi(\mu))
\quad\text{and}\quad X(y,0)=y,\quad\Psi(0)=0
\]
\Unclear{GELFAND-P000474}{middle-induced-variables}{The first arguments of $G$ and $X$ are heavily retraced and provisionally read as $x$. The right side clearly has $X(y,0)=y$.}


\uline{Rule for constructing a vers. family:}
Take an orbit and a set transverse to the orbit.
\SourceFigure{GELFAND-P000474}{transverse-sketch}{A curved line and a short transverse segment with a point at their intersection.\par\smallskip\centering
[DIAGRAM]}

{\centering\rule{1.8em}{.4pt}\quad\textperiodcentered\quad\rule{2.4em}{.4pt}\par}

\[
F(x,\lambda)=f(x)+
\sum a_i\left.\frac{\partial F}{\partial\lambda_i}\right|_{\lambda_i=0}
=f(x)+\sum c_i\varphi_i(x)+O(\lambda^2).
\]
\Unclear{GELFAND-P000474}{lower-taylor-coefficient}{The multiplier in the first sum is provisionally read as $a_i$. The evaluation subscript is $\lambda_i=0$.}


\uuline{\mbox{$\varphi_i$} give}\quad\uline{an infinitesimal versal deformation:}
if $\forall$ fn. $a(x)$ is repr. in the form
\Unclear{GELFAND-P000474}{lower-definition-heading}{The first fragment is heavily retraced and double-underlined. The provisional reading is $\varphi_i$ «зад» (“give”).}

\[
a(x)=\sum\underline{\underline{\frac{\partial f}{\partial x_i}}}\cdot h_i(x)
+\sum_{i=1}^{k}c_i\varphi_i(x)
\]
$a(x)$ defines an infinitesim. vers. deform.

($\varphi_i$ gen. the local ring of smooth maps)

\Incomplete{GELFAND-P000474}{bottom-cancelled-taylor}{A separate crossed-out entry precedes the parenthesis about Taylor expansion. Its reading is unresolved.}

(from the Taylor exp. of that fn.)

Important: the quotient ring by the \uline{Jacobian ideal}
$\left(\left\{\dfrac{\partial f}{\partial x_i}\right\}\right)$.

\uline{The quotient} — a finite-dimensional space of dimension $k$.

{\centering\rule{2.9em}{.4pt}\quad\textperiodcentered\quad\rule{3.4em}{.4pt}\par}


\Needspace{7\baselineskip}

\SourcePage{GELFAND-P000475}{unlabelled}
\[
\underline{\underline{f=z_1^{k+1}+z_2^2+\cdots+z_n^2.}}
\]

Then
\[
\frac{\partial f}{\partial z_1}=z_1^k,\qquad
\frac{\partial f}{\partial z_2}=z_2,\ \ldots,
\]
\Unclear{GELFAND-P000475}{top-derivatives}{The right-hand sides are written as monomials without numerical factors.}

hence the quotient ring is a space of dimension $k$.

\[
\varphi_i=1,z_i,z_i^2,\ldots,z_1^{k-1}
\]
\Unclear{GELFAND-P000475}{top-phi-indices}{The first two powers after $1$ have subscripts resembling $i$. The last has $1$.}


Vers. deformation:
\[
\underline{\lambda_0+\lambda_1z_1+\cdots+\lambda_{k-1}z_1^{k-1}
+z_1^{k+1}+z^2}
\quad\text{— \uline{vers. deformation}}
\]
\SourceFigure{GELFAND-P000475}{versal-bracket}{A rule with upturned ends is drawn below the polynomial. The final term is written $z^2$.}


{\centering\rule{2.8em}{.4pt}\quad\textperiodcentered\quad\rule{3.2em}{.4pt}\par}
\uline{\mbox{$F(z,\lambda)=0$}.}\qquad
The bifurc. diagram — the set of those $\lambda$, for which sing. points remain (in the vers. deformation).

\SourceFigure{GELFAND-P000475}{theorem-star}{Four segments meet at a common point to the left of the theorem.\par\smallskip\centering
[DIAGRAM]}

\uuline{Thm.:} on an open
\(\underset{\text{\scriptsize\tikz[baseline=(label.base)]\node[draw,ellipse,inner sep=1pt] (label) {bifurc. diag.};}}{\text{set}}\)
the sing. points will already be simple (Morse).
\SourceFigure{GELFAND-P000475}{theorem-bifurcation-insertion}{“Bifurc. diag.” is written in a bubble below “set”.}


\[
z_1^{k+1}+z_2^2=0
\qquad\Vert\qquad
\begin{gathered}
\text{the bifurc. diag. of this point —}\\[-.2ex]
\longrightarrow\text{ Swallowtail}.
\end{gathered}
\]
\[
\begin{gathered}
z_1^{k+1}+z_2^2+\lambda_{k-1}z_1^{k-1}+\cdots+\lambda_0\\
z_2^2=-p_{k+1}(z_1);\qquad z_2=\sqrt{-p_{k+1}(z_1)}.
\end{gathered}
\]
sing. points\quad $(z_2,P'(z_1))=0$, i.e.
\[
\underline{P(z_1)=0\ \text{and}\ P'(z_1)=0}
\]
($P$ has multiple roots)
\SourceFigure{GELFAND-P000475}{multiple-root-bracket}{A single rule with upturned ends surrounds the two equalities below.}


Consider the complement of the bifurc. diagram and its fund. group.
It is the braid group of the Coxeter group.

Consider the eq.
\[
\underline{P(z,\lambda)=0,\quad |z|<\varepsilon}
\qquad\text{\uline{\mbox{$\lambda$} such that there are no singularities}}
\]
\Unclear{GELFAND-P000475}{middle-polynomial-name}{The first letter is heavily retraced. The provisional reading is $P$.}


the homology of this space is $\ne0$ only in \sout{dim.}
\Incomplete{GELFAND-P000475}{lower-homology-cancelled}{Traces remain after the struck-out “dim.”. A full reading is unresolved.}

half-dimension $n-1$.

It is a bouquet of $\mu$ spheres.
\SourceFigure{GELFAND-P000475}{sphere-bouquet}{Four loops with a common point depict the bouquet.\par\smallskip\centering
[DIAGRAM]}


\Incomplete{GELFAND-P000475}{lower-left-cancelled}{The entry to the left of $H^{n-1}$ is heavily struck out. Its reading is unresolved.}

$H^{n-1}$ has generators $e_1\ e_2\ldots\ e_\mu$.

The intersection index defines a scalar prod. $(e_i,e_k)$.

\[
n-1\equiv2\pmod4\qquad\Vert\qquad
\begin{gathered}
\text{neg. def.}\\
\text{semidef.}\\
\text{indef.}
\end{gathered}
\qquad\underline{e_i^2=-2}\quad\text{(always)}
\]
\Incomplete{GELFAND-P000475}{lower-signature-cancelled}{One line between “semidef.” and “indef.” is crossed out. Its reading is unresolved.}


\uwave{monodromy}: rep. of the fund. gr. in
\Incomplete{GELFAND-P000475}{lower-monodromy-cancelled}{An entry after “in” is crossed out and followed by an ellipsis. Its reading is unresolved.}

\ldots

{\centering\rule{4em}{1pt}\quad\textbullet\quad\rule{4em}{1pt}\par}

\sout{Atiyah}
\Unclear{GELFAND-P000475}{bottom-cancelled-name}{The struck-out name may read «Атья» (“Atiyah”).}

\uline{There are works by Atiyah, Bott,\ldots:} how to get rid of derivatives in Singer's calculations?
How can one study the topology of the world while sitting at one point? Finding explicit asymptotics for eqs. of the 2nd order (considerations of I. M. Gelfand, V. I. Arnold and others)

(will be in the next semester)

A. Shnirelman on Maslov's work.


\Needspace{7\baselineskip}

\SourcePage{GELFAND-P000476}{unlabelled}
\uuline{V. M. Buchstaber.} (continuation).

\[
\begin{gathered}
A\overset{\text{addit.}}{=}\Lambda[[X]],\qquad
X_n\text{ — addit. generators}\\
X_n\cdot X_m=C_{n+m}^n X_{n+m}
\end{gathered}
\]
\Unclear{GELFAND-P000476}{additive-equality}{Two short parallel strokes under “addit.” read as an equals sign.}


\[
\begin{gathered}
\operatorname{Ext}_A^{2,*}(\Lambda,\Lambda)\\[-.4ex]
\rotatebox[origin=c]{90}{$=$}\\[-.4ex]
\operatorname{Ker}d_3^*/\operatorname{Im}d_2^*
\end{gathered}
\qquad\text{in terms of the }B\text{-construction}
\]

$1^\circ$.
\[
d_3^*\colon\overline A^*\otimes\overline A^*
\longrightarrow\overline A^*\otimes\overline A^*\otimes\overline A^*
\]

\[
A^*=\operatorname{Hom}_{\Lambda}^{c}(\Lambda[[X]],\Lambda)
\overset{\text{add.}}{\simeq}\Lambda[x]
\]
\Unclear{GELFAND-P000476}{top-dual-hom}{The upper index of $\operatorname{Hom}$ reads $c$ and the lower one $\Lambda$.}

\uuline{compact} (nonzero val. on a finite number).
\SourceFigure{GELFAND-P000476}{compact-annotation}{“Compact” is written below $\operatorname{Hom}$ and connected to it by a stroke.}


\[
\begin{gathered}
\Delta\colon A^*\longrightarrow A^*\otimes A^*\\
\Delta x_n\longrightarrow x_n\otimes1+1\otimes x_n+\sum d_{i,j}x_i\otimes x_j^*
\end{gathered}
\]
\Unclear{GELFAND-P000476}{top-delta-star}{A star is visible after $x_j$. It is absent from the later formula for $\Delta x_n$.}


\[
\begin{aligned}
(\Delta x_n,x_i\otimes x_j)&=d_{i,j}\\[-.4ex]
\qquad\rotatebox[origin=c]{90}{$=$}&\\[-.4ex]
(x_n,x_i\cdot x_j)&=(x_n,C_{i+j}^{i}x_{i+j})=C_{i+j}^{i}\delta_{n,i+j}
\end{aligned}
\]
\SourceFigure{GELFAND-P000476}{pairing-layout}{A vertical equals sign joins the two lines of the pairing.}


\[
\Delta x_n=x_n\otimes1+1\otimes x_n+\sum C_n^i x_i\otimes x_{n-i}
\]
And put\quad $x_n=x^n$
\[
\begin{gathered}
\Delta x^n=(x\otimes1+1\otimes x)^n\\
d_3^*\colon \overline{\Lambda[x,y]}\longrightarrow\overline{\Lambda[x,y,z]}
\end{gathered}
\]

\[
d_3\Biggl(\underset{
  \substack{\rotatebox[origin=c]{90}{$=$}\\ \sum h_{i,j}x^iy^j}
}{h(x,y)}\Biggr)=\sum h_{ij}(\overline\Delta x^i)z^j
-\sum h_{i,j}x^i\overline\Delta y^j=
\]
\Unclear{GELFAND-P000476}{middle-differential-star}{The calculation has $d_3$ without a star before $h(x,y)$. The preceding line has $d_3^*$.}


\[
\overline\Delta x^n=
\underset{\substack{
 \rotatebox[origin=c]{90}{$=$}\\
 (\Delta(x))^n\\
 \rotatebox[origin=c]{90}{$=$}\\
 (x+y)^n
}}{\Delta x^n}-x^n-y^n.
\]
\SourceFigure{GELFAND-P000476}{coproduct-identifications}{Below $\Delta x^n$, vertical equals signs connect $(\Delta(x))^n$ and $(x+y)^n$.}


\[
=h(x+y,z)-\mbox{\sout{\ensuremath{h(x,z)}}}-h(y,z)
-h(x,y+z)+h(x,y)+\mbox{\sout{\ensuremath{h(x,z)}}}.
\]
\SourceFigure{GELFAND-P000476}{cancelled-cocycle-terms}{Both instances of $h(x,z)$ are crossed out. The minus and plus signs in front remain visible.}


\uuline{\mbox{$\operatorname{Ker}d_3^*$}:}
\[
\boxed{h(x+y,z)+h(x,y)=h(x,y+z)+h(y,z)}
\]
Denote these cycles by $L$.
\SourceFigure{GELFAND-P000476}{kernel-relation-box}{The equality is enclosed in a rectangular box.}


This is precisely the Lazard relation.
\Unclear{GELFAND-P000476}{lazard-name}{The adjective before “relation” reads «лазаровское» (“Lazard”).}


\[
\underline{\operatorname{Im}d_2^*}=\Lambda\{d_2^*x^n\}
\]
\[
0\longrightarrow\operatorname{Im}d_2^*\longrightarrow L
\longrightarrow\operatorname{Ext}_A^{2,n}(\Lambda,\Lambda)\longrightarrow0
\]

\uline{Calculation of the group \mbox{$\operatorname{Ext}$}.}

The eq. does not involve $\Lambda$. Hence it suffices to calc. 2 cases:

1)\quad $\Lambda$ — $Z_p$-module ($p$ — prime)

2)\quad $\Lambda=Q$


\Needspace{7\baselineskip}

\SourcePage{GELFAND-P000490}{unlabelled}
\uline{Milnor–Kervaire thm..}\quad All $M_1$ are obtained from one another by Morse surgeries.

(necess. and suff.)

Is there a manifold there homot. equiv. to $M$.

$P_n$ — the obstruction group
\Unclear{GELFAND-P000490}{obstruction-cancellation}{A word after “obstructions” is crossed out. Its reading is unresolved.}

for surgery to a homotop. equiv.
\[
\underline{\sigma_0:M_1\longrightarrow M}
\]
\uline{The seq. is exact.}
\[
?\longrightarrow hS(M^n).
\]
\SourceFigure{GELFAND-P000490}{exact-question}{A question mark is written in place of the first term.}


\[
\begin{aligned}
\alpha_1.&\quad\{N_1\overset{f_1}{\longrightarrow}M\},\\
\alpha_2:&\quad\{N_2\overset{f_2}{\longrightarrow}M\}.
\end{aligned}
\]
\SourceFigure{GELFAND-P000490}{two-structures}{A period follows $\alpha_1$, and a colon follows $\alpha_2$.}


$\xi_1\in\phi(M)$ is constr. from $\alpha_1$.
\SourceFigure{GELFAND-P000490}{group-sign}{The sign before $(M)$ is set as $\phi$.}


\SourceFigure{GELFAND-P000490}{thom-pair}{Two arrows point to Thom complexes. Two vertical strokes stand between the targets.\par\smallskip
{\centering
[DIAGRAM]\par}}

\Unclear{GELFAND-P000490}{bundle-symbol}{The wavy stroke may be part of $\tau_M$ or a separate tilde over $\tau$. It is provisionally set as $\tau_M$.}


We obtain the problem:
\Unclear{GELFAND-P000490}{novikov-cancellation}{An entry between “problem” and “diss.” is crossed out. Its reading is unresolved.}

Novikov’s diss.:
\SourceFigure{GELFAND-P000490}{homotopy-target}{Below the target are $\cup$ and $M$.\par\smallskip
{\centering
[DIAGRAM]\par}}


Film:
\[
\begin{aligned}
F\big|_{S^L\times0}&=\beta_1\\
F\big|_{S^L\times1}&=\beta_2
\end{aligned}
\]
\SourceFigure{GELFAND-P000490}{endpoint-restrictions}{The two restrictions of $F$ at $0$ and $1$.}

\Unclear{GELFAND-P000490}{last-line-cancellation}{The sign before the quantifier is overwritten. Its reading is unresolved.}


$\exists$  $h$-cobordisms among the films?


\Needspace{7\baselineskip}

\SourcePage{GELFAND-P000491}{unlabelled}
\uline{19/\mbox{$\overline{\underline{\mathrm{III}}}$}-73}
\qquad\uline{Gelfand’s Sem.}
\Unclear{GELFAND-P000491}{date-cancellation}{A short entry before the date is crossed out. It may contain digits.}


\uuline{Arnold.}\quad Consid. the germ of a fn. $f:(\mathbb C^n,0)\longrightarrow(\mathbb C,0)$.
$0$ — a crit. point.
\SourceFigure{GELFAND-P000491}{annotation-icon}{A yellow callout icon with three black strokes is visible above the line. Its text is unavailable in this frame.}


Consid. the space of $k$-jets. The group of $k$-jets of diffs. acts.

$j^k f$\qquad $k$-suff. (the subsequent ones have no effect)

\uline{\mbox{$m$} — number of moduli.}
\uuline{Thm..}\quad $m=0\Longrightarrow$ the list
\[
\begin{array}{ll}
A_k & x^{k+1}\\
D_k & x^2y+y^{k-1}\\
E_6 & x^3+y^4\\
E_7 & x^3+xy^3\\
E_8 & x^3+y^5
\end{array}
\]
\SourceFigure{GELFAND-P000491}{ade-table}{Several strokes in the list are retraced, including those in $E_6$ and $E_8$.}


Stab. equiv. \quad $\oplus$ with nondeg. quadr. forms.

The quadr. form is positive definite.


\Needspace{7\baselineskip}

\SourcePage{GELFAND-P000492}{unlabelled}
An analyt. fn., up to an analyt. change of var., uniquely det. the change of var.

$0$-modality sing. — after changing the coeffs. $\exists$ only a finite number of variants.
\Unclear{GELFAND-P000492}{modality-ending}{The ending was overwritten. The line is provisionally read as «модульная», with «модульных» also possible.}


{\centering\rule{1.4em}{.4pt}\quad\textperiodcentered\quad\rule{1.6em}{.4pt}\par}

$f=\varepsilon$\quad level line:
\uline{Intersection with a ball} — an even-dimensional manifold with boundary.
\[
f^{-1}(\varepsilon)\cap B_\delta
\]
\SourceFigure{GELFAND-P000492}{fibre-drawing}{Two closed curves are joined by inward-curving sides.\par\smallskip
{\centering
[DIAGRAM]\par}}


Intersection index
\Unclear{GELFAND-P000492}{intersection-cancellation}{An entry after “intersection” is crossed out. Its reading is unresolved.}

of real homology in the middle dim. (integral)

It is
$\overset{\text{\scriptsize (neg. def.)}}{\text{definite in sign}}$
exactly in the listed cases,
provided $n\equiv3\ (\mathrm{mod}\ 4)$.
(otherwise add extra variables) — Tyurin’s thm..
\SourceFigure{GELFAND-P000492}{negative-insert}{“(Neg. def.)” is inserted in a callout after “definite in sign”.}

\[
f^{-1}(\varepsilon)\cap B_\delta\approx
\underbrace{S^{n-1}\vee\ldots\vee S^{n-1}}_{\mu}.
\]
$\mu=\dim$ of the local ring.

\SourceFigure{GELFAND-P000492}{dynkin-diagram}{Diagram of $E_6$. The leftmost vertex is marked $-2$ and the first edge is marked $1$.\par\smallskip
{\centering
[DIAGRAM]\par}}

\Unclear{GELFAND-P000492}{graph-cancellation}{An entry to the right of $E_6$ is crossed out. Its reading is unresolved.}


The quadr. form is given by the Dynkin diag.

\fbox{$m=1$.} — one also finds a list (long).

First:
[DIAGRAM]

\begin{center}
[DIAGRAM]
\end{center}

[DIAGRAM]
$\Longrightarrow$ always $m=1$.
\Unclear{GELFAND-P000492}{modality-inequality}{The sign before $1$ is retraced. It is provisionally set as $\leqslant$. The conclusion reads $m=1$.}


For all of them
\[
\mu_+\leqslant1\Longrightarrow\mu_0=1\qquad(\text{always degen.})
\]
\uuline{\mbox{$\mu_+=1$}.}\quad
$axyz+x^k+y^\ell+z^m$ — \uline{three-index series} $T_{k\ell m}$.

\uuline{A. S. Mishchenko.}

\uline{Sullivan:} localization, periodicity and \ldots\ Galois-
\[
K(X)/\text{\mbox{$J$}-equiv.}=J(X)\quad\text{— \uline{a finite group}}.
\]
Adams: it is described in terms of the cohom. ops. of $K$-theory $\psi^k$.
\[
\psi^k(\xi)=\xi^k\quad\text{and \uwave{additive}}
\quad\Longrightarrow\quad\text{this determines them uniquely}
\]
$(\dim\xi=1)$\qquad $k=2,3,\ldots$

Thus $J(X)$ is determined by the group $K(X)$ as a module over these ops..

$J$-equiv. — fibrewise homot. equiv..
\SourceFigure{GELFAND-P000492}{fibrewise-diagram}{The arrow $f$ joins $\xi$ and $\eta$. Both letters have arrows to $X$.\par\smallskip
{\centering
[DIAGRAM]\par}}

$f$ — proper. \quad $\deg f=1$ on fibres.

It was proved (Adams, Atiyah):
\[
\xi\overset{J}{\sim}\text{triv.}.
\]

$X$ — a finite complex.
\raisebox{.45ex}{\scriptsize (then $\exists N\in\mathbb Z$:)}
\SourceFigure{GELFAND-P000492}{integer-quantifier}{“Then $\exists N\in\mathbb Z$:” is inserted after “complex”.}


$\forall$ $J$-triv. bundles $\xi$
\raisebox{.45ex}{\scriptsize ($\forall k$)}
$\exists$
\SourceFigure{GELFAND-P000492}{k-quantifier}{The quantifier $\forall k$ is inserted between $\xi$ and $\exists$.}

\Unclear{GELFAND-P000492}{bundle-cancellation}{A word before “bundle” is crossed out. Its reading is unresolved.}

a bundle $y$:
\Unclear{GELFAND-P000492}{formula-cancellation}{An entry after the colon is crossed out. Its reading is unresolved.}

\begin{center}
[DIAGRAM]
\end{center}
\SourceFigure{GELFAND-P000492}{adams-equality-bracket}{A stroke runs below and to either side of the equality.}


\uline{The converse — the Adams conjecture:}
$\forall$ finite compl. $X$ $\exists N$; $\forall k$
\[
k^N(\psi^k-1)\xi\overset{J}{\sim}0.
\]


\Needspace{7\baselineskip}

\SourcePage{GELFAND-P000493}{unlabelled}
\uline{Idea.}
\SourceFigure{GELFAND-P000493}{universal-bundle-diagram}{The left side reads $\xi=f^*(\xi)$ and the right side $\xi_n$.\par\smallskip
{\centering
[DIAGRAM]\par}}

\Unclear{GELFAND-P000493}{pullback-star}{The sign after $f$ is heavily retraced and provisionally set as a star.}


\SourceFigure{GELFAND-P000493}{pullback-square}{The second diagram shows $f^*(\eta)$ and $\eta$.\par\smallskip
{\centering
[DIAGRAM]\par}}

\[
\psi^k(f^*(\eta))=f^*\psi^k(\eta).
\]
Hence the op. $\psi^k$ is det. by specifying a bundle on \uuline{\mbox{$BU(n)$}}.
$(\psi^k(\xi_n))$.

From $\psi^k$ one constructs a map
$BU(n)\longrightarrow BU(\infty)$ (unknown dimension).

$\psi^k(1)=1\Longrightarrow$ it follows that $\exists$ a sequence of maps.
\SourceFigure{GELFAND-P000493}{stabilization-diagram}{Above $BU(n)$ is $\xi_n+1$.\par\smallskip
{\centering
[DIAGRAM]\par}}


Thus $\psi^k$ is given by a single map.
\begin{center}
[DIAGRAM]
\end{center}
\SourceFigure{GELFAND-P000493}{single-arrow}{A double stroke runs below the arrow, with turned-up ends.}


\uwave{Suppose there is a cohom. op.}
$BU(\infty)\overset{A}{\longrightarrow}BU(\infty)$,
and suppose it preserves dim.
\SourceFigure{GELFAND-P000493}{rank-square}{The square has $A$ on both horizontal arrows.\par\smallskip
{\centering
[DIAGRAM]\par}}


\uwave{preserves dimension.}
\uwave{They do not occur, but if they did, one could prove that}
$A\xi\overset{J}{\sim}\xi$.
\SourceFigure{GELFAND-P000493}{dimension-margin}{The note on the right has a wavy underline.}


Then $\underline{A\xi\overset{J}{\sim}\xi}$.

[DIAGRAM]
\SourceFigure{GELFAND-P000493}{spheres-label}{The frame around “spheres” is open on the left.}

\SourceFigure{GELFAND-P000493}{sphere-square}{The diagram has two levels and vertical ellipses.\par\smallskip
{\centering
[DIAGRAM]\par}}

\SourceFigure{GELFAND-P000493}{left-base-notation}{The lower-left node reads $B(n+1)$.}

\Unclear{GELFAND-P000493}{right-sphere-arrow}{The upper-right arrow is heavily retraced. The arrowhead is visible on the left.}


A fibrewise homot. equiv. is obtained.

\sout{But: how to}
\SourceFigure{GELFAND-P000493}{question-cancellation}{“But: how to” is crossed out.}

There are none, but how do we “find” them?
\SourceFigure{GELFAND-P000493}{j-factorization}{Two arrows go from $X$ to $BU(n)$.\par\smallskip
{\centering
[DIAGRAM]\par}}

(\uline{classif. space of bundles with fibre a}
$\overset{\text{\scriptsize (homotopy)}}{\text{\uline{sphere}}}$)
\SourceFigure{GELFAND-P000493}{homotopy-sphere-insert}{“Homotopy” is inserted before “sphere”.}


$J$ — a functor.

Suppose $\exists BU(n)^{\wedge}$, to which the Adams operations transfer,
i.e. $\exists$ an op. $(\psi^k)^{\wedge}$,
for which the diag. commutes.
\SourceFigure{GELFAND-P000493}{completed-adams-diagram}{The lower-right arrow bends up toward $B_F$.\par\smallskip
{\centering
[DIAGRAM]\par}}

Perhaps, it may be that $(\psi^k)^{\wedge}$
\uline{preserves dimension.}

(\uline{orientations} — dim.-preserving maps)

These spaces were constr. by Sullivan.


\Needspace{7\baselineskip}

\SourcePage{GELFAND-P000494}{unlabelled}
$BU(n)^{\wedge}$ will be the compl. Grassmannian.

On
\Unclear{GELFAND-P000494}{initial-cancellation}{A short entry after “On” is crossed out. Its reading is unresolved.}

\uuline{\mbox{$S^{2n-1}$} the group \mbox{$U(n)$}}.
\Unclear{GELFAND-P000494}{sphere-exponent}{The sign in the exponent is heavily retraced. It is provisionally set as $2n-1$, possibly $2n+1$.}

We obtain bundles (spher.) on $BU(n)$.

\uline{Sullivan:} another space $\widehat S^{2n-1}$ and the group $G_n$.
Consid. the bundle and the classif. space.
\SourceFigure{GELFAND-P000494}{top-completion-diagram}{A stroke runs below the horizontal arrow. The separate entry on the right is $\widehat S^{2n+1}$.\par\smallskip
{\centering
[DIAGRAM]\par}}

and the fibrewise homot. type will be preserved.

\Incomplete{GELFAND-P000494}{obliterated-line}{The line is heavily crossed out. A continuous reading is unresolved.}


Simply connected compls. with fin. gen. homot. grps. are considered.

\uline{Claim.}\quad
$\overset{\forall X}{\exists}$
\SourceFigure{GELFAND-P000494}{forall-insert}{$\forall X$ is inserted above $\exists$.}

\SourceFigure{GELFAND-P000494}{universal-triangle}{A stroke runs under the upper arrow. Two arrows point to the bottom node.\par\smallskip
{\centering
[DIAGRAM]\par}}

(functor)

$\alpha(X)$ — a transformation of the ident. functor.

($\forall F$, \raisebox{.5ex}{\scriptsize (for which)}
$\pi_i(F)$ are finite).
\SourceFigure{GELFAND-P000494}{finite-insert}{“For which” is inserted after $F$.}

\Unclear{GELFAND-P000494}{finite-cancellation}{An entry before “are finite” is blacked out. Its reading is unresolved.}


\uuline{\mbox{$\pi_i(\widehat X)$} are profinite}
($\underset{\longleftarrow}{\lim}$ fin. groups).
(in partic., compact).

\sout{In particular}
\Unclear{GELFAND-P000494}{before-then-cancellation}{The crossed-out words are provisionally read as “In particular”.}

Then one can compute:
\[
\pi_i(\widehat X)=\widehat{\pi_i(X)}
\]
(profin. completion).
Neighb. — subgroups of fin. index and their limit.

$Z^{\wedge}$ is descr. by
$\overset{\text{\scriptsize (integral)}}{\text{\mbox{$p$}-adic numbers}}$.
\SourceFigure{GELFAND-P000494}{padic-insert}{“Integral” is inserted before $p$.}


The homot. groups of $\widehat X$ behave like integral $p$-ad. numbers.

\sout{Action of the Galois grp..}
\SourceFigure{GELFAND-P000494}{galois-cancellation}{The heading is crossed out.}

\SourceFigure{GELFAND-P000494}{completion-square}{The lower arrow points right and has an equals sign above it.\par\smallskip
{\centering
[DIAGRAM]\par}}

all homot. grps. are finite (these are stable homot. groups of spheres)

\SourceFigure{GELFAND-P000494}{lower-completion-diagram}{The lower horizontal arrow points right. There is an additional stroke at its left end. The arrow from $\widehat S^{2n+1}$ crosses the neighbouring writing.\par\smallskip
{\centering
[DIAGRAM]\par}}


\sout{Suppose}
\Unclear{GELFAND-P000494}{adams-cancellation}{The following short word is crossed out. Its reading is unresolved.}

\raisebox{.5ex}{\scriptsize (The constr. op. $\psi_k$)}
extends the Adams ops. and preserves dimension.
\Unclear{GELFAND-P000494}{adams-operation-index}{The small sign to the right of $\psi$ is provisionally set as subscript $k$. Another possible reading is $\widehat\psi^k$.}


Consid. the Grassmannians: $G_{N,n}$.
There are finitely many aff. alg. charts, with eqs. over $Z$ (scheme over $Z$).
— an alg. set over $Q$.
$\Rightarrow$ the Galois grp. acts there $\operatorname{Gal}(\mathbb C:Q)$
— acts on $G_{N,n}$.
It is discontinuous in the usual sense, but continuous in the Zariski top..
Consid. \uuline{Étale} top. \ldots
\[
\operatorname{Gal}(\mathbb C:Q)\longrightarrow
\operatorname{Gal}\bigl(
\underset{\substack{\vert\\\text{\scriptsize max. ab. ext.}}}{Q_{ab}};Q\bigr)
=\widehat Z^{*}
\]
\SourceFigure{GELFAND-P000494}{abelian-label}{Below $Q_{ab}$ is the label “max. ab. ext.”.}

\[
H_i(V_c,A)\approx H_i(V_{et},A)
\]
\[
\widehat V_c\sim\underset{\longleftarrow}{\lim}\,N(V_{et}),
\qquad N\text{ — nerve}.
\]


\Needspace{7\baselineskip}

\SourcePage{GELFAND-P000495}{unlabelled}
\uuline{26.3.73}

\uwave{Statement of the res. of} \uwave{\mbox{\uuline{Atiyah, Patodi, Singer.}}} \uwave{(I. M. Gelfand).}

Let $A$ — be a self-adj. ell. op.
\[
\zeta_A(s)=\operatorname{Tr}A^{-s}=\sum\lambda_i^{-s}
\]
at $s=0$ has no pole (Seeley).
\Unclear{GELFAND-P000495}{seeley-spelling}{The Cyrillic spelling of the name is provisionally read as «Сили», possibly «Силли».}


$\zeta_A(0)$ — Atiyah and others.

Consid. In place of the $\zeta$-fn., $L$-series.
\Incomplete{GELFAND-P000495}{eta-cancelled-phrase}{The crossed-out phrase to the right may begin “Not”. The rest is illegible.}


\begin{flushright}\small \uline{Let \mbox{$A$} — not pos. def.} (self-adj. ellip.)\end{flushright}
\[
\eta(s)=\sum_{\lambda\ne0}(\operatorname{sign}\lambda_i)\,|\lambda_i|^{-s}
\]
Difficult theor.: \uline{\mbox{$s=0$} — is no pole} !!

Important example — 1st-order op.

General prob.: no pos. def. Consid.
$\sum$ (mult. characters of $\lambda_i$)
($L$-series).
\uline{On compl. numbers} \uline{there are more complicated mult. characters.}

\uline{Example.}
\Unclear{GELFAND-P000495}{example-cancellation}{The crossed-out entry before $A$ partly overlaps $A$.}

\[
 A=i\frac{d}{dx}+t\qquad\text{on }[0,2\pi]\qquad\text{periodic}
\]
\[
\underline{\eta(0)=1-2t}
\]
\SourceFigure{GELFAND-P000495}{eta-value-domain}{An open-topped bracket runs under the formula.}


Eigenvals. $n+t$, $n=-2,-1,0,1,\ldots$

$\eta(0)$ — the “asymmetry” of the operator:
\[
\sum_{\lambda_i>0}1-\sum_{\lambda_i<0}1.
\]

\uline{S. I. Gelfand.} Account of Manin’s last lecture:

$p$-adic integration (Mazur, Swinnerton — Dyer)
\Unclear{GELFAND-P000495}{swinnerton-spelling}{The middle letters of the first part of the name run together. Provisionally read as Swinnerton.}



\Needspace{7\baselineskip}

\SourcePage{GELFAND-P000496}{unlabelled}
$Z_p$ — integral $p$-adic numbers, a fn. is defined on them. It is to be integrated.

($p$-adic vals.)

\uuline{lit.} L. G. Schnirelmann, Dwork, Grauert, Krasner, Tate

$Z_p\supset A_i=p^iZ_p$ — syst. of neighb. of $0$.
The measure is def. on sets $a+A_i$

$\mu(a+A_i)$ — a $p$-adic number (not nec. integral).
\[
a\in Z_p\quad\text{or}\quad \underline{a\in Z_p/p^i}
\]
There must be finite additivity; reduces to the cond.:
\[
\mu(a+A_i)=\sum_{t=0}^{p-1}\mu(a+tp^i+A_{i+1})
\]

\uline{2 examples}: 1) \uline{invar. measure}: indep. of $a$, but dep. on $i$.
Then
\[
\mu(a+A_i)=p^{-i}\qquad\text{(only possibility).}
\]
The measure of a small set is large! — far from $0$
\[
\left(\left\|\frac ab\right\|=p^{-n},\quad
\text{where }\frac ab=p^n\frac{a_1}{b_1},\quad
a_1,b_1\text{ rel. prime to }p\right).
\]

2) \uline{More important example.} (Leopoldt — Kubota measures).
\begin{flushright}
2 parameters $n$ and $a$
\end{flushright}
\[
(n,p)=1,\qquad 0\le a\le p^i-1
\]
\[
\mu_n(a+A_i)=\left[\frac{na}{p^i}\right]-\frac{n-1}{2}
\]
\[
\mu_2(a+A_i)=
\begin{cases}
+1/2,&\text{if }a>\dfrac{p^i}{2}\\
-1/2,&\text{if }a<\dfrac{p^i}{2}
\end{cases}
\]
It is addit. !! The measure of $\forall$ set is bounded

(it is obtained by the Mellin transf. from the $\vert$ $\zeta$-fn. of Leopoldt — Iwasawa)
\Unclear{GELFAND-P000496}{mellin-vertical-stroke}{A vertical stroke appears between “from” and $\zeta$. Its function is unresolved.}


\[
Q_p\ni f(x),\quad x\in Z_p\qquad
\int_{Z_p}f(x)\mu(x)\qquad\text{in Riemann’s sense}
\]
$R_m$ — representatives of $Z_p/p^mZ_p$ (1 from $\forall$ class)

This is a partition. Form
\[
S(R_m)=\sum_{a\in R_m}f(a)\mu(a+A_m)
\]
they converge under conds. on the measure and the fn.

\uline{Moderately growing measure} — if it grows more slowly than the invariant one
\[
\lim_{m\to\infty}\max|\mu(a+A_m)|p^{-m}=0
\]
$f$ satisfies the Lipschitz cond. if
\[
 |f(a)-f(b)|<Cp^{-m},\quad\text{if }|b-a|<p^{-m}
\]
(in partic., if given by a conv. power series)

If the measure grows moderately and the fn. satisfies the Lipschitz cond., then, $\exists$ a limit, called the integral.


\Needspace{7\baselineskip}
